% Data and Methodology slides
\begin{frame}{Dữ liệu}
  \textbf{Dataset:} \texttt{WA\_Fn-UseC\_-HR-Employee-Attrition.csv} với 1{,}470 quan sát
  
  \vspace{0.4cm}
  \textbf{Biến phân loại chính:}
  \begin{itemize}
    \item OverTime, JobRole, BusinessTravel, MaritalStatus
  \end{itemize}
  
  \vspace{0.4cm}
  \textbf{Biến số tiêu biểu:}
  \begin{itemize}
    \item Age, MonthlyIncome, YearsAtCompany
    \item EnvironmentSatisfaction, JobInvolvement, DistanceFromHome
  \end{itemize}
  
  \vspace{0.4cm}
  \textbf{Biến mục tiêu:} \texttt{Attrition} (1 = Nghỉ việc, 0 = Ở lại)
\end{frame}

\begin{frame}{Dữ liệu - Lưu ý}
  \begin{alertblock}{Lưu ý}
    Tỷ lệ nghỉ việc 16.12\% $\Rightarrow$ mất cân bằng lớp (class imbalance) $\Rightarrow$ dùng \texttt{class\_weight='balanced'}
  \end{alertblock}
  
  \vspace{0.4cm}
  \textbf{Ý nghĩa:}
  \begin{itemize}
    \item Dataset chỉ có khoảng 16\% nhân viên nghỉ việc thực tế
    \item Thuật toán LR dễ bị thiên hướng luôn dự đoán "Không nghỉ việc"
    \item \texttt{class\_weight='balanced'} giúp mô hình học nhạy cảm hơn đối lưu các phần tử thiểu số
  \end{itemize}
\end{frame}

\begin{frame}{Logistic Regression - Bài toán và Công thức}
  \textbf{Bài toán:} Tính xác suất một nhân viên sẽ đưa ra quyết định nghỉ việc (Attrition = 1).
  
  \vspace{0.2cm}
  \textbf{Công thức xác suất Sigmoid:}
  \[
  p(y=1\mid x) = \frac{1}{1+e^{-z}}, \quad \text{với } z = \beta_0 + \beta_1 x_1 + \dots + \beta_d x_d
  \]
  
  \textbf{Tính toán từng phần tử để thay vào công thức:}
  \begin{itemize}
    \item \textbf{$p(y=1\mid x)$:} Kết quả đầu ra (xác suất làm tròn nội suy từ 0 đến 1).
    \item \textbf{Biến hệ $x_i$:} Giá trị đặc trưng của nhân viên (VD: $x_1$ = Age, $x_2$ = OverTime=Yes).
    \item \textbf{Trọng số $\beta_i$:} Mức độ tác động của đặc trưng (Do thuật toán tính toán).
    \item \textbf{Tính $z$:} Nhân thực tế từng giá trị $x_i$ của nhân viên đó với hệ số $\beta_i$, rồi lấy tổng cộng với hệ số tự do $\beta_0$.
    \item \textbf{Toán hạng hàm mũ $e$:} Thế $z$ vào $1+e^{-z}$ ở mẫu số để ánh xạ kết quả $z$ tuyến tính thành một xác suất toán học.
  \end{itemize}
\end{frame}

\begin{frame}{Hàm mục tiêu và Regularization}
  \textbf{Hàm mục tiêu:} Âm log-likelihood (NLL) kết hợp L2 regularization 
  \[
  \mathcal{L}(\beta) = -\sum_{i=1}^{n}\Big[y_i\log p_i + (1-y_i)\log(1-p_i)\Big] + \lambda \|\beta\|_2^2
  \]
  
  \vspace{0.4cm}
  \begin{itemize}
    \item Regularization ngăn hệ số mô hình bùng nổ, dễ xử lý các biến số đo lường đơn vị lớn.
  \end{itemize}
  
  \vspace{0.3cm}
  \begin{alertblock}{Thiết kế Split \& Tiền Xử Lý}
    \begin{itemize}
      \item Train/Test theo stratified 80/20 nhằm đảm bảo phân bố 16\% xuyên suốt.
      \item sklearn: Pipeline (OneHotEncoder + StandardScaler) phục vụ hệ thống máy học tốc độ cao.
      \item statsmodels: Cung cấp p-value và tính Odds Ratio giải thích chuyên sâu.
    \end{itemize}
  \end{alertblock}
\end{frame}
